\subsection{Существующие подходы и архитектуры}

Существует ряд подходов и архитектур, которые решают данную задачу. Рассмотрим основные из них.
%-------------------------


\begin{enumerate}
    \item Generative Adversarial Networks (GAN)

    Генеративные состязательные сети (GAN) стали основой для множества задач обработки изображений и видео. Их успех связан с использованием двух взаимодействующих сетей — генератора и дискриминатора. Генератор создает реалистичные данные, а дискриминатор оценивает их подлинность, улучшая способность генератора создавать правдоподобные результаты.

    \begin{itemize}
        \item Видео-ориентированные GAN: Расширенные версии GAN, такие как MoCoGAN и TGAN, адаптированы для обработки видеоданных, учитывая временные зависимости между кадрами. Эти архитектуры используют рекуррентные модули для моделирования последовательностей, что делает их пригодными для задач, где необходимо сохранять временную согласованность.
        
        \item Стилизация и изменение формы объектов: GAN активно применяются для стилизации видео, создания глубоких фейков и замены объектов. Например, StyleGAN и его модификации обеспечивают высокую степень контроля над стилем и деталями объектов в кадре.

    \end{itemize}

    Однако GAN имеют ограничения, связанные с трудностями в обучении и необходимостью большого объема данных для генерации качественных результатов.

    \item Рекуррентные нейронные сети (RNN) и их модификации

    Рекуррентные нейронные сети хорошо зарекомендовали себя для обработки последовательных данных, включая видеопотоки. Благодаря своим механизмам памяти они способны учитывать временную информацию.

    \begin{itemize}
        \item Long Short-Term Memory (LSTM): LSTM-архитектуры широко применяются для трекинга объектов и анализа последовательностей. Они позволяют эффективно обрабатывать длительные временные зависимости, что критически важно для задач динамической замены объектов.

        \item Gated Recurrent Units (GRU): GRU предлагают упрощенную структуру по сравнению с LSTM, обеспечивая при этом схожую производительность. Они используются в задачах, где важна эффективность вычислений.
    \end{itemize}

    \item Диффузионные модели
    
    Диффузионные модели представляют собой относительно новое направление, которое стремительно развивается. Эти модели основаны на постепенном добавлении и удалении шума из данных, что позволяет им генерировать изображения и видео высокой реалистичности.

    \begin{itemize}
        \item FateZero: Современная методология, представленная в статье FateZero, использует диффузионные модели для редактирования видеоданных. Ключевой особенностью подхода является использование карт внимания (“attention maps”) на этапах инверсии и денойзинга для сохранения временной согласованности.

        \item Пространственно-временные модификации: Пространственно-временные механизмы внимания (spatial-temporal attention) позволяют моделям учитывать как пространственные, так и временные аспекты видео. Это критически важно для обеспечения естественного движения объектов после их модификации.
        \item Генеративные модели с нулевым обучением (zero-shot): FateZero демонстрирует способность к редактированию видео без дополнительного обучения на конкретных задачах, что делает этот подход особенно перспективным для широкого применения.
    \end{itemize}

    \item Подходы на основе трансформеров

    Трансформеры, изначально разработанные для обработки текстов, показали свою эффективность в задачах обработки изображений и видео. Их способность обрабатывать длинные последовательности делает их особенно подходящими для задач с временными зависимостями.

    \begin{itemize}
        \item Vision Transformers (ViT): ViT используется для сегментации объектов и анализа сцен. Это позволяет выделять ключевые области видео для замены или модификации объектов.

        \item Трансформеры для видео (Video Transformers): Адаптированные версии трансформеров, такие как TimeSformer, используют механизм внимания для анализа временных зависимостей между кадрами, сохраняя высокую точность обработки.
    \end{itemize}

    \item Оптический поток

    Методы, основанные на анализе оптического потока, позволяют отслеживать движение объектов и учитывать их динамику при замене или модификации. Это классический подход, который часто используется в сочетании с глубокими моделями для повышения их эффективности.

    \begin{itemize}
        \item FlowNet: Один из первых успешных подходов для оценки оптического потока, который позволяет моделям понимать движение объектов в сцене.

        \item RAFT: Современная архитектура, обеспечивающая высокую точность оценки оптического потока, используется для сложных сцен с большим количеством движущихся объектов.
    \end{itemize}

    \item Гибридные подходы
    
    Для достижения наилучших результатов современные исследования стремятся объединить несколько методов. Например, комбинация диффузионных моделей, GAN и трансформеров позволяет получить более качественные результаты за счет использования сильных сторон каждой технологии.

    \begin{itemize}
        \item Интеграция GAN и трансформеров: Сочетание генеративных возможностей GAN с мощными механизмами внимания трансформеров обеспечивает как реалистичность, так и точность замены объектов.

        \item Диффузионные модели и оптический поток: Объединение диффузионных моделей с методами оценки оптического потока позволяет учитывать движение объектов, сохраняя их согласованность в видеопотоке.
    \end{itemize}
\end{enumerate}

Современные архитектуры для динамической замены и модификации объектов в видеоданных демонстрируют значительные успехи благодаря интеграции передовых технологий, таких как GAN, трансформеры и диффузионные модели. Исследования в этой области продолжаются, и новые подходы, такие как FateZero, открывают возможности для более эффективного и реалистичного редактирования видеоданных. В следующей главе будут рассмотрены практические аспекты применения описанных архитектур, а также разработка предложенного алгоритма для решения поставленной задачи.

%------------


